% chapter: 計算の基本

\chapter{計算の基本}\label{章：計算の基本}
\begin{multicols*}{2}
  漸化式の解説を行う前に,実際の計算において基本となる内容をこの章では説明する.
  ここで扱う内容は漸化式の一般項を得る過程だけでなく,様々な場面での計算において有用なものである.
  ごく狭い場面で使えるものは避け,多くの場面で使われる手法と,それらに通底する思想を取り扱っている.
  前半では和の計算,式変形,指数・対数を確認し,続いて数学的帰納法と二項定理を扱う.
  後半では三角関数,極形式,複素数の順に,数列の値を角度や平面上の点として表すための知識を準備する.
  複素数の節では,先に説明した三角関数の加法定理と数学的帰納法を用いる.
% section: シグマの計算

\section{シグマの計算}\label{節：シグマの計算}
  \noindent\textbf{【本編で使う場面】}\par
  和の記号と添字の役割を確認する.\sectionref*{節：階差型}で変化を足し戻し,\sectionref*{節：総和型}で総和の式をずらして引くときに使う.

  多くの項を足すとき、すべてを並べる代わりに、一つの記号でまとめられる。
  実数の\bookterm{sequence}{数列}$\{a_k\}$について、$a_1$から$a_n$までの和を
  \[
    \sum_{k=1}^{n} a_k=a_1+a_2+\cdots+ a_n
  \]
  と書く。複素数の\bookterm{sequence}{数列}でも、同じ記号を使う。
  $\sum$は、ギリシャ文字の\bookterm{sigma}{シグマ}である。
  下の$k=1$から上の$n$まで、$k$を一つずつ進めて足す。

  同じ数を繰り返し足せば、項数を掛けた値になる。
  定数倍は和の外へ出せる。二つの和は、項ごとに分けてもよい。
  \begin{align*}
    &\sum_{k=1}^{n} c=nc\\
    &\sum_{k=1}^{n} ca_k=c\sum_{k=1}^{n} a_k\\
    &\sum_{k=1}^{n} (a_k+b_k)=\sum_{k=1}^{n} a_k+\sum_{k=1}^{n} b_k
  \end{align*}
  自然数と、その二乗・三乗の和は、次の式でまとめられる。
  \begin{align*}
    &\sum_{k=1}^{n} k=\frac{1}{2}n(n+1)\\
    &\sum_{k=1}^{n} k^2=\frac{1}{6}n(n+1)(2n+1)\\
    &\sum_{k=1}^{n} k^3=\left\{\frac{1}{2}n(n+1)\right\}^2\\
  \end{align*}

  添字$k$は,和の中で順に動く番号である.同じ範囲なら$\sum_{k=1}^{n}a_k=\sum_{j=1}^{n}a_j$と書き換えても値は変わらない.
  一方,上端の$n$は足す項数を決める.例えば$S_n=\sum_{k=1}^{n}a_k$なら
  \[
    S_{n+1}=S_n+a_{n+1}
  \]
  である.
  \bookterm{recurrence}{漸化式}の添字をずらす場合は,式全体の$n$を置き換える.
  例えば$a_{n+1}=2a_n+n$から$a_{n+2}=2a_{n+1}+n+1$を得る.

% section: 式の変形

\section{式の変形}\label{節：式の変形}
  \noindent\textbf{【本編で使う場面】}\par
  同じ値を保ちながら既知の形を作る操作を確認する.定数の調整は\sectionref*{節：特性方程式型},\bookterm{partial}{部分分数分解}は\sectionref*{節：階差型}で和を計算する際に使う.

  \subsection{\texorpdfstring{$+a-a$をする}{+a-aをする}}\label{節：+a-aをする}
    二次式の中に、二乗の形を作る。
    $a\ne0$として、足りない定数を足し、同じだけ引くと、
    \begin{align*}
      &ax^2+bx+c\\
      &\qquad=a\left(x^2+\frac{b}{a}x\right) +c\\
      &\qquad=a\left(x^2+\frac{b}{a}x+\frac{b^2}{4a^2}-\frac{b^2}{4a^2}\right)+c \tag*{\(\cdots\) (☆)}\\
      &\qquad=a\left(x+\frac{b}{2a}\right)^2-\frac{b^2}{4a}+c
    \end{align*}
    となる。☆の行で、二乗に必要な定数を補った。
    足して引いた量は零なので、式の値は変わらない。
    $x$を含む部分が一つの二乗にまとまり、式の形が見やすくなる。
    この操作を\bookterm{square}{平方完成}という。
    二次関数の最大・最小や、軸の位置を調べるときにも使える。

    値を保ったまま、扱いやすい形を作る。
    同じ発想が、後の定数項を取り除く操作につながる。

  \subsection{部分分数分解}\label{節：部分分数分解}
    次は、積を含む分母を、より簡単な分数の和へ分ける。
    \[
      \frac{1}{(x+a)(x+b)}
    \]
    は、そのままでは和を計算しにくい。
    分母を一次の式に分けるため、次の形を作る。
    ここでは$x+a\ne0$、$x+b\ne0$とする。
    \[
      \frac{1}{(x+a)(x+b)}=\frac{\alpha}{x+a}+\frac{\beta}{x+b}
    \]
    $\alpha,\beta$について求めると,
    \begin{align*}
    \frac{1}{(x+a)(x+b)}
      &= \frac{\alpha}{x+a}+\frac{\beta}{x+b}\\
      &= \frac{(\alpha+\beta)x+\alpha b+\beta a}
      {(x+a)(x+b)}
    \end{align*}
    \[
    \therefore\quad
    \begin{cases}
    \alpha+\beta=0,\\
    \alpha b+\beta a=1
    \end{cases}
    \]
    これを解くと,\,$\beta=-\alpha$であるから,
    \[
      \alpha b-\alpha a=1
      \iff \alpha(b-a)=1
    \]
    したがって,\,$a\ne b$であれば,
    \[
      \alpha=\frac{1}{b-a},
      \qquad
      \beta=-\frac{1}{b-a}
    \]
    が得られる(\,$a=b$のときはそもそも分母が重解$(x+a)^2$になり,この形の分解はできない).
    二つの係数を定めると、一次の分母をもつ分数の和に分かれる。
    この操作を\bookterm{partial}{部分分数分解}という。
    3次の場合も,\,$p,q,r\ne0$で,分母の一次因子$px+a,qx+b,rx+c$の根
    \[
      -\frac{a}{p},\qquad -\frac{b}{q},\qquad -\frac{c}{r}
    \]
    が互いに異なるときは,
\BookMathLayout{\[
      \frac{1}{(px+a)(qx+b)(rx+c)}=\frac{\alpha}{px+a}+\frac{\beta}{qx+b}+\frac{\gamma}{rx+c}
    \]}{
\[
\begin{aligned}
\frac{1}{(px+a)(qx+b)(rx+c)}
&=\frac{\alpha}{px+a}+\frac{\beta}{qx+b}\\
&\quad+\frac{\gamma}{rx+c}
\end{aligned}
\]
}
    と置き,係数$\alpha,\beta,\gamma$を定めることで分解できる.
    言い換えれば,これらの一次因子が互いに定数倍でないことが条件である.
    重複因子があるときは,その重複度に応じて例えば$\dfrac{A}{x+a}+\dfrac{B}{(x+a)^2}$のような項も必要になる.

% section: 等比数列の総和とその導出

\section{等比数列の総和とその導出}\label{節：等比数列の総和とその導出}
  初項が$a$,公比が$r$である等比数列
  \[
    a,\,ar,\,ar^2,\,\ldots,\,ar^{n-1}
  \]
  の初項から第$n$項までの総和を考える.
  その総和を$S_n$とすると,
  \[
    S_n
    =\sum_{k=1}^{n} ar^{k-1}
    =a+ar+ar^2+\cdots+ar^{n-1}
  \]
  と表せる.

  \subsection{総和の公式}
    公比$r$が$1$でないとき,
    \[
      S_n=\frac{a(1-r^n)}{1-r}
    \]
    である.
    分母と分子の符号を同時に変えれば,
    \[
      S_n=\frac{a(r^n-1)}{r-1}
    \]
    とも書ける.
    公比が$1$のときは,すべての項が$a$であるから,
    \[
      S_n=na
    \]
    となる.

  \subsection{公式の導出}
    $r\ne1$として,総和の式の両辺に$r$を掛ける.
    \begin{align*}
      S_n
        &=a+ar+ar^2+\cdots+ar^{n-1},\\
      rS_n
        &=ar+ar^2+ar^3+\cdots+ar^n.
    \end{align*}
    2つの式を項ごとに引くと,右辺では共通する項がすべて消える.
    \begin{align*}
      S_n-rS_n
        &=(a+ar+ar^2+\cdots+ar^{n-1})\\
        &\qquad -(ar+ar^2+ar^3+\cdots+ar^n)\\
        &=a-ar^n.
    \end{align*}
    よって,
    \[
      (1-r)S_n=a(1-r^n)
    \]
    である.
    $r\ne1$より両辺を$1-r$で割ると,
    \[
      S_n=\frac{a(1-r^n)}{1-r}
    \]
    を得る.

    この導出では,総和を公比$r$倍してから元の総和との差を取ることで,途中の項を消している.
    数列の和を求めるときには,このように総和を定数倍して引き,共通する項を消す考え方がよく用いられる.


% section: 指数対数

  \section{指数対数}
    \subsection{指数}
      $a$が$p$回掛けられている状況を考える.
      ここで
      \[
        a\times a\times\cdots\times a
      \]
      のように記すのは煩雑なので,
      \[
        a\times a\times\cdots\times a=a^p
      \]
      と表す.
      このときの$p$を指数といい,
      \[
        f(x)=a^x
      \]
      で表されるような関数を指数関数という.
      ここでは$p$を正の整数として導入したが,\,$a^0=1,\ a^{-p}=\dfrac1{a^p}$のように定義を拡張すれば,指数が$0$や負の整数,有理数,実数であっても同様に$a^x$を考えることができる.
      以降,特に断らない限りこの拡張された意味で指数を用いる.
    \subsection{対数}
      $a^p=q$のとき,\,$p$は
      \[
        \log_a{q}=p
      \]
      のように表される.
      この時$\log$をログと読み,\,$p$を対数という.
      また,\,$a$を底,\,$q$を真数と言う.
      一般に底$a$は$a>0,\,a\ne1$を満たす必要がある.
      また,
      \[
        e=\lim_{x \to \infty} \left(1+\frac{1}{x}\right)^x
      \]
      で定義される数$e$を自然対数の底またはネイピア数という.また,
      \[
        f(x)=\log_a x
      \]
      というような関数を対数関数といい,指数関数の逆関数である.
      逆関数とは
      \[
        y=f(x)
      \]
      を満たすような$x,y$について,
      \[
        x=g(y)
      \]
      となるような関数$g$を指す.
      グラフ上では,逆関数は元の関数と関数$y=x$に対して対称である.
      関数$g$はよく,\,$f^{-1}$と書かれる.

      対数に関して,一般に以下が成り立つ.
      \begin{align*}
        &\log_a{p^q}=q\log_a{p}\\
        &\log_a{MN}=\log_a{M}+\log_a{N}\\
        &\log_a{\frac{M}{N}}=\log_a{M}-\log_a{N}\\
        &\log_a{b}=\frac{\log_c{b}}{\log_c{a}}
      \end{align*}
    

% section: 数学的帰納法

  \section{数学的帰納法}
    自然数$n$についての主張$P(n)$が,すべての$n\ge1$で成り立つことを示したいとする.
    このとき,以下の2つを示せばよい.
    \begin{enumerate}[label=(\roman*),parsep=3pt]
      \item $P(1)$が成り立つ.
      \item ある$n$で$P(n)$が成り立つと仮定すると,\,$P(n+1)$も成り立つ.
    \end{enumerate}
    この2つが示されれば,\,(i)より$P(1)$が成り立ち,\,(ii)より$P(1)\Rightarrow P(2)$,\,$P(2)\Rightarrow P(3)$,\,$\ldots$と次々に成り立つことが分かるので,すべての$n\ge1$について$P(n)$が成り立つ.
    ドミノ倒しに例えられることが多い.\,(i)は最初のドミノを倒すこと,\,(ii)は「あるドミノが倒れれば次のドミノも倒れる」ように並べることに対応する.

    本書では,漸化式から予想した一般項の形を確かめる場面や,三角関数の加法定理から漸化式を導く場面(3.2や3.6.3など)で数学的帰納法を用いる.
% section: 二項定理

  \section{二項定理}
    次に,和の累乗を展開するための公式を確認する.
    自然数$n$と実数$x,y$について,
    \[
      (x+y)^n=\sum_{k=0}^{n}\binom{n}{k}x^{n-k}y^{k}
    \]
    が成り立つ.これを二項定理という.ここで$\dbinom{n}{k}$は二項係数であり,
    \[
      \binom{n}{k}=\frac{n!}{k!(n-k)!}
    \]
    で定義される. ここで,$k$は$0\le k\le n$を満たす整数であり,$0!=1$とする.
    この二項係数は,$n$個から$k$個を選ぶ組合せの数に一致する.
    例えば$n=2$では,
    \[
      (x+y)^2=x^2+2xy+y^2
    \]
    $n=3$では,
    \[
      (x+y)^3=x^3+3x^2y+3xy^2+y^3
    \]
    となる.
    後の節で扱う複素数に対しても,同じ二項定理が成り立つ.
    付録では,この公式を用いて差分を取ると多項式の次数が下がる理由などを説明する.
\section{三角関数}
  ここからは,数列の値を角度や平面上の点として表すための準備を行う.
  まず,極形式や複素数の計算でも用いる三角関数の公式を振り返る.
  まず一般に,
  \begin{align*}
    \sin(\alpha+\beta)
      &=\sin\alpha\cos\beta+\cos\alpha\sin\beta\\
    \cos(\alpha+\beta)
      &=\cos\alpha\cos\beta-\sin\alpha\sin\beta\\
    \tan(\alpha+\beta)
      &=\frac{\tan\alpha+\tan\beta}
              {1-\tan\alpha\tan\beta}
  \end{align*}
  が成り立つ.
  ここで $\alpha=\beta=x$ とすれば,倍角公式
  \begin{align*}
    \sin 2x&=2\sin x\cos x\\
    \cos 2x&=2\cos^2x-1=1-2\sin^2x\\
    \tan 2x&=\frac{2\tan x}{1-\tan^2x}
  \end{align*}
  を得る.また,倍角公式から
  \begin{align*}
    \sin^2x&=\frac{1-\cos 2x}{2}\\
    \cos^2x&=\frac{1+\cos 2x}{2}\\
    \tan^2x&=\frac{1-\cos 2x}{1+\cos 2x}
  \end{align*}
  を導くことができる.また,加法定理から,\,$3x=2x+x$とすることで,
  \begin{align*}
    \sin 3x&=3\sin x-4\sin^3x\\
    \cos 3x&=4\cos^3x-3\cos x
  \end{align*}
  が求められる.
\section{極形式}
  \subsection{極座標系の導入}
    前節で確認した三角関数を用いると,平面上の点を距離と角度によって表すことができる.
    普段用いている直交座標系と比較しながら,この表し方を導入する.
    直交座標系とは,デカルト座標系とも呼ばれ,平面上で直交する2直線の交点を基準として任意の点を表す手法のことである.
    しかしながら円や媒介変数表示された関数などにおいて,一部煩雑な表現をせざるをえない状況があった.
    また,図形の平行移動や対称移動に関しては容易に行えたものの,回転させたり歪ませたりするような操作は非常に複雑な操作が必要だった.
    そこで,直交座標系に加え,新たな概念として極座標系の導入を行う.
    
    極座標系とは,平面上の任意の点を原点からの距離と角度を用いて記述する方法である.
    例えば

    \begin{tikzpicture}[scale=0.8]
      % 座標軸
      \draw[->,>=stealth,semithick] (-1.5,0)--(5.5,0) node[above]{$x$};
      \draw[->,>=stealth,semithick] (0,-1)--(0,5) node[right]{$y$};

      % 原点
      \draw (0,0) node[below right]{O};

      % 点 (1,sqrt(3)) への補助線
      \draw[dashed] (1,0) -- (1,{sqrt(3)});
      \draw[dashed] (0,{sqrt(3)}) -- (1,{sqrt(3)});

      % 座標軸上の値
      \draw (1,0) node[below]{$1$};
      \draw (0,{sqrt(3)}) node[left]{$\sqrt{3}$};

      % 点
      \fill (1,{sqrt(3)}) circle (2pt)
        node[above right]{$(1,\sqrt{3})$};
    \end{tikzpicture}

    のように表される点は,極座標系においては

\begin{tikzpicture}[scale=0.8]

  % 座標軸
  \draw[->,>=stealth,semithick]
    (-1.5,0)--(5.5,0) node[above]{$x$};
  \draw[->,>=stealth,semithick]
    (0,-1)--(0,5) node[right]{$y$};

  % 原点・点A
  \coordinate (O) at (0,0);
  \coordinate (A) at (1,{sqrt(3)});
  \node[below left] at (O) {O};

  % 動径
  \draw[very thick] (O)--(A);

  % 下側の曲線
  \draw[semithick]
    (O)
    .. controls (0.02,0.35) and (0.12,0.62)
    .. (0.3,0.7);

  % 上側の曲線
  % 終点Aでの接線を動径OAと平行にする
  \draw[semithick]
    (0.6,1.4)
    .. controls (0.62,1.58) and
               (0.82,{sqrt(3)-0.10*sqrt(3)})
    .. (A);

  % 動径の長さ
  \node at (0.30,1.05) {$2$};

  % 角度
  \draw[semithick]
    (0.55,0) arc (0:60:0.55);
  \node at (0.68,0.3) {$\frac{\pi}{3}$};

  % 点
  \fill (A) circle (1.5pt);

\end{tikzpicture}
    と表される. この点の原点からの距離は$2$,正の$x$軸から測った角度は$\pi/3$である.
  \subsection{極形式}
    一般に,原点からの距離を$\rho$,正の$x$軸から反時計回りに測った角度を$\theta$とすると,
    点の直交座標$(x,y)$は
    \[
      x=\rho\cos\theta,\qquad y=\rho\sin\theta
    \]
    と表される. ここで$\rho=\sqrt{x^2+y^2}$である.
    原点以外では$\rho>0$であり,$\theta$に$2\pi$の整数倍を加えても同じ点を表す.
    原点では$\rho=0$となり,角度は一意に定まらない.

% section: 複素数の極形式と指数表示

\section{複素数の極形式と指数表示}

前節では,平面上の点を原点からの距離と角度によって表した.
ここでは,その点を複素数として読み替える.
$i^2=-1$を満たす虚数単位$i$を用いると,実数$x,y$に対して複素数は$z=x+iy$と表される.
この複素数を平面上の点$(x,y)$に対応させたものを複素平面という.
前節の$x=\rho\cos\theta,\,y=\rho\sin\theta$を代入すれば,
\[
  z
  =
  \rho(
    \cos\theta+i\sin\theta
  )
\]
と書くことができる. $\rho$は複素数の絶対値であり,原点からの距離を表す.
$z\ne0$のとき,$\theta$を偏角という. 偏角は$2\pi$の整数倍の違いを除いて定まる.

このとき,
\[
  \cos\theta+i\sin\theta
\]
のような複素数を何乗かしたときに,
偏角がどのように変化するかを考える.

まず,三角関数の加法定理を用いる.


\subsection{極形式の積}

  任意の実数$\alpha,\beta$について,
  \begin{align*}
    &(\cos\alpha+i\sin\alpha)
    (\cos\beta+i\sin\beta)\\
    &=
    \cos\alpha\cos\beta
    -\sin\alpha\sin\beta\\
    &\qquad
    +i(
      \sin\alpha\cos\beta
      +\cos\alpha\sin\beta
    ).
  \end{align*}

  加法定理より,
  \[
    \cos\alpha\cos\beta
    -\sin\alpha\sin\beta
    =
    \cos(\alpha+\beta)
  \]
  および
  \[
    \sin\alpha\cos\beta
    +\cos\alpha\sin\beta
    =
    \sin(\alpha+\beta)
  \]
  であるから,
  \[
    (\cos\alpha+i\sin\alpha)
    (\cos\beta+i\sin\beta)
    =
    \cos(\alpha+\beta)
    +i\sin(\alpha+\beta)
  \]
  を得る.

  つまり,
  \[
    \cos\alpha+i\sin\alpha
  \]
  と
  \[
    \cos\beta+i\sin\beta
  \]
  を掛けると,絶対値は1のままで偏角が
  \[
    \alpha+\beta
  \]
  となる.

  この性質を繰り返し用いることで,
  複素数の累乗について重要な公式が得られる.


\subsection{ド・モアブルの定理}

  任意の整数$n$について,
  \[
    (\cos\theta+i\sin\theta)^n
    =
    \cos(n\theta)+i\sin(n\theta)
  \]
  が成り立つ.

  まず,$n$が正の整数の場合を考える.

  $n=1$のとき,
  \[
    (\cos\theta+i\sin\theta)
    =
    \cos\theta+i\sin\theta
  \]
  であるから,成り立つ.

  いま,
  \[
    (\cos\theta+i\sin\theta)^n
    =
    \cos(n\theta)+i\sin(n\theta)
  \]
  が成り立つと仮定する.

  このとき,
  \begin{align*}
    &(\cos\theta+i\sin\theta)^{n+1}\\
    &=
    \{
      \cos(n\theta)+i\sin(n\theta)
    \}
    (\cos\theta+i\sin\theta)\\
    &=
    \cos((n+1)\theta)
    +i\sin((n+1)\theta)
  \end{align*}
  となる.

  よって,数学的帰納法により,
  \[
    (\cos\theta+i\sin\theta)^n
    =
    \cos(n\theta)+i\sin(n\theta)
  \]
  がすべての正の整数$n$について成り立つ.

  $n=0$のときは,
  \[
    (\cos\theta+i\sin\theta)^0
    =
    1
    =
    \cos0+i\sin0
  \]
  である.

  また,
  \[
    \cos(-\theta)=\cos\theta,
    \qquad
    \sin(-\theta)=-\sin\theta
  \]
  より,
  \begin{align*}
    &(\cos\theta+i\sin\theta)
    (\cos(-\theta)+i\sin(-\theta))\\
    &=
    \cos0+i\sin0\\
    &=
    1.
  \end{align*}

  したがって,
  \[
    (\cos\theta+i\sin\theta)^{-1}
    =
    \cos(-\theta)+i\sin(-\theta)
  \]
  である.

  これを用いれば,負の整数についても同じ公式が成り立つ.

  したがって,
  \[
    (\cos\theta+i\sin\theta)^n
    =
    \cos(n\theta)+i\sin(n\theta)
    \qquad(n\in\mathbb Z)
  \]
  が成り立つ.

  これをド・モアブルの定理という.

  例えば,
  \[
    \left(
      \cos\frac{\pi}{6}
      +i\sin\frac{\pi}{6}
    \right)^6
    =
    \cos\pi+i\sin\pi
    =
    -1
  \]
  となる.

  このように,複素数を極形式で表すと,累乗を考えるときには絶対値の累乗と偏角の$n$倍を考えればよいことが分かる.


\subsection{複素数の指数表示}

  次に,実数の指数関数で用いられる
  \[
    e^x
    =
    \lim_{n\to\infty}
    \left(
      1+\frac{x}{n}
    \right)^n
  \]
  という極限表示を複素数に拡張する.

  ここでは,
  \[
    e^{i\theta}
    :=
    \lim_{n\to\infty}
    \left(
      1+\frac{i\theta}{n}
    \right)^n
  \]
  と定める.

  この極限を評価するため,
  \[
    1+\frac{i\theta}{n}
  \]
  を極形式に直す.

  その絶対値は
  \[
    r_n
    =
    \sqrt{
      1+\frac{\theta^2}{n^2}
    }
  \]
  である.

  また,$\phi_n$を
  \[
    \tan\phi_n
    =
    \frac{\theta}{n}
  \]
  を満たす角とする.

  このとき,
  \[
    \cos\phi_n
    =
    \frac{1}{r_n}
  \]
  および
  \[
    \sin\phi_n
    =
    \frac{\theta}{nr_n}
  \]
  であるから,
  \[
    1+\frac{i\theta}{n}
    =
    r_n(
      \cos\phi_n+i\sin\phi_n
    )
  \]
  と表すことができる.

  よって,
  \begin{align*}
    \left(
      1+\frac{i\theta}{n}
    \right)^n
    &=
    r_n^n
    (\cos\phi_n+i\sin\phi_n)^n\\
    &=
    r_n^n
    \{
      \cos(n\phi_n)
      +i\sin(n\phi_n)
    \}.
  \end{align*}

  最後の等式にはド・モアブルの定理を用いた.

  したがって,この極限を求めるためには,
  \[
    r_n^n
  \]
  と
  \[
    n\phi_n
  \]
  の極限を調べればよい.


\subsection{極限の計算}

  まず,
  \[
    r_n^n
    =
    \left(
      1+\frac{\theta^2}{n^2}
    \right)^{n/2}
  \]
  について考える.

  自然対数の基本極限
  \[
    \lim_{x\to0}
    \frac{\log(1+x)}{x}
    =
    1
  \]
  を用いると,
  \begin{align*}
    \log(r_n^n)
    &=
    \frac n2
    \log
    \left(
      1+\frac{\theta^2}{n^2}
    \right)\\
    &=
    \frac{\theta^2}{2n}
    \frac{
      \log
      \left(
        1+\frac{\theta^2}{n^2}
      \right)
    }{
      \frac{\theta^2}{n^2}
    }.
  \end{align*}

  $n\to\infty$のとき,
  \[
    \frac{\theta^2}{2n}\to0
  \]
  であり,
  \[
    \frac{
      \log
      \left(
        1+\frac{\theta^2}{n^2}
      \right)
    }{
      \frac{\theta^2}{n^2}
    }
    \to1
  \]
  であるから,
  \[
    \log(r_n^n)\to0.
  \]

  よって,
  \[
    r_n^n\to1
  \]
  である.

  次に,
  \[
    n\phi_n
  \]
  の極限を考える.

  $\tan\phi_n=\dfrac{\theta}{n}$より,
  \[
    \tan\phi_n\to0
  \]
  であるから,
  \[
    \phi_n\to0
  \]
  である.

  また,
  \[
    \tan x
    =
    \frac{\sin x}{\cos x}
  \]
  および,
  \[
    \lim_{x\to0}\frac{\sin x}{x}=1,
    \qquad
    \lim_{x\to0}\cos x=1
  \]
  より,
  \[
    \lim_{x\to0}\frac{\tan x}{x}=1
  \]
  である.

  したがって,
  \[
    \lim_{x\to0}\frac{x}{\tan x}=1
  \]
  であるから,
  \begin{align*}
    n\phi_n
    &=
    \theta
    \frac{\phi_n}{\tan\phi_n}\\
    &\to
    \theta.
  \end{align*}

  よって,
  \[
    n\phi_n\to\theta
  \]
  が成り立つ.


\subsection{オイラーの公式}

  以上より,
  \[
    r_n^n\to1
  \]
  および
  \[
    n\phi_n\to\theta
  \]
  である.

  $\sin x,\cos x$は連続関数であるから,
  \[
    \cos(n\phi_n)\to\cos\theta
  \]
  および
  \[
    \sin(n\phi_n)\to\sin\theta
  \]
  となる.

  したがって,
  \begin{align*}
    e^{i\theta}
    &=
    \lim_{n\to\infty}
    \left(
      1+\frac{i\theta}{n}
    \right)^n\\
    &=
    \lim_{n\to\infty}
    r_n^n
    \{
      \cos(n\phi_n)
      +i\sin(n\phi_n)
    \}\\
    &=
    \cos\theta+i\sin\theta.
  \end{align*}

  よって,
  \[
    e^{i\theta}
    =
    \cos\theta+i\sin\theta
  \]
  が成り立つ.

  これをオイラーの公式という.

  したがって,
  \[
    \cos\theta+i\sin\theta
    =
    e^{i\theta}
  \]
  と書くことができる.

  これにより,複素数
  \[
    z
    =
    \rho(
      \cos\theta+i\sin\theta
    )
  \]
  は,
  \[
    z
    =
    \rho e^{i\theta}
  \]
  と表せる.


\subsection{ド・モアブルの定理の指数表示}

  オイラーの公式を用いると,
  ド・モアブルの定理は指数関数を使ってさらに簡潔に表すことができる.

  実際,
  \[
    \cos\theta+i\sin\theta
    =
    e^{i\theta}
  \]
  より,
  \begin{align*}
    (\cos\theta+i\sin\theta)^n
    &=
    (e^{i\theta})^n\\
    &=
    e^{in\theta}\\
    &=
    \cos(n\theta)+i\sin(n\theta).
  \end{align*}

  また,
  \[
    z
    =
    \rho e^{i\theta}
  \]
  と表せば,
  \begin{align*}
    z^n
    &=
    \rho^n e^{in\theta}\\
    &=
    \rho^n
    \{
      \cos(n\theta)+i\sin(n\theta)
    \}.
  \end{align*}

  したがって,
  \[
    z^n
    =
    \rho^n
    \{
      \cos(n\theta)+i\sin(n\theta)
    \}
    \qquad(n\in\mathbb Z)
  \]
  となる.

  ここでは,極形式における「絶対値は$n$乗し,偏角は$n$倍する」という性質が,
  指数表示によって簡潔に表されている.


\subsection{実数指数への拡張}

  これまでのド・モアブルの定理では,指数$n$は整数であった.

  ここで,指数を一般の実数にまで広げることを考える.

  まず,複素数
  \[
    z
    =
    \rho(
      \cos\theta+i\sin\theta
    )
    =
    \rho e^{i\theta}
    \qquad(\rho>0)
  \]
  に対して,実数$x$を指数とする場合を考える.

  実数の指数関数では,
  \[
    \rho^x
  \]
  が定義されているので,
  \[
    z^x
  \]
  を
  \[
    z^x
    :=
    \rho^x e^{ix\theta}
  \]
  と定める.

  オイラーの公式より,
  \[
    e^{ix\theta}
    =
    \cos(x\theta)+i\sin(x\theta)
  \]
  であるから,
  \[
    z^x
    =
    \rho^x
    \{
      \cos(x\theta)+i\sin(x\theta)
    \}
  \]
  となる.

  したがって,
  \[
    z^x
    =
    \rho^x
    \{
      \cos(x\theta)+i\sin(x\theta)
    \}
    \qquad(x\in\mathbb R)
  \]
  と表すことができる.

  $x$が整数の場合には,
  これはすでに示したド・モアブルの定理と一致する.

  また,$x=\dfrac{m}{n}$が有理数の場合には,
  \[
    z^{m/n}
    =
    \rho^{m/n}
    \left\{
      \cos\frac{m\theta}{n}
      +i\sin\frac{m\theta}{n}
    \right\}
  \]
  である.

  この数を$n$乗すると,
  \begin{align*}
    \left(
      z^{m/n}
    \right)^n
    &=
    \rho^m
    \left(
      \cos\frac{m\theta}{n}
      +i\sin\frac{m\theta}{n}
    \right)^n\\
    &=
    \rho^m
    \{
      \cos(m\theta)
      +i\sin(m\theta)
    \}\\
    &=
    z^m.
  \end{align*}

  したがって,$z^{m/n}$は$z^m$の$n$乗根となっており,
  この定義は有理数の場合にも通常の指数の意味と整合している.

  以上から,ド・モアブルの定理は,
  \[
    z^x
    =
    \rho^x
    \{
      \cos(x\theta)+i\sin(x\theta)
    \}
    \qquad(x\in\mathbb R)
  \]
  という形に拡張することができる.

  ただし,複素数の偏角は一般には一意ではなく,
  \[
    \theta,\quad
    \theta+2\pi,\quad
    \theta+4\pi,\quad\ldots
  \]
  はすべて同じ複素数を表す.

  このため,一般の複素数の実数乗を考えるときには,
  どの偏角を用いるかをあらかじめ固定しておく必要がある.

  本書では,
  \[
    z
    =
    \rho(
      \cos\theta+i\sin\theta
    )
  \]
  において$\theta$を一つ固定し,
  \[
    z^x
    =
    \rho^x
    \{
      \cos(x\theta)+i\sin(x\theta)
    \}
  \]
  と考えることにする.
\end{multicols*}
